\documentclass[UTF8,zihao=-4]{ctexart}
\usepackage[a4paper,margin=2.4cm]{geometry}
\usepackage{fontspec}
\usepackage{xcolor}
\usepackage{booktabs}
\usepackage{longtable}
\usepackage{array}
\usepackage{calc}
\usepackage{graphicx}
\usepackage{hyperref}
\usepackage{enumitem}
\usepackage{fancyhdr}
\usepackage{tikz}
\usetikzlibrary{arrows.meta,positioning,fit,calc,shapes.geometric}

\setCJKmainfont{Microsoft YaHei}
\setmainfont{Times New Roman}
\hypersetup{colorlinks=true,linkcolor=blue!50!black,urlcolor=blue!60!black}
\providecommand{\tightlist}{\setlength{\itemsep}{0pt}\setlength{\parskip}{0pt}}
\setlist[itemize]{leftmargin=2em}
\setlist[enumerate]{leftmargin=2em}
\pagestyle{fancy}
\setlength{\headheight}{15pt}
\fancyhf{}
\lhead{MEMO-Appflow}
\rhead{端侧 eBPF + MAPLE 实时调度}
\cfoot{\thepage}

\title{\textbf{MEMO-Appflow：基于真实 eBPF 证据与端侧大模型的 Android 实时应用预测和系统调度}}
\author{Jingyi Guo}
\date{2026-07-01}

\begin{document}
\maketitle

\begin{abstract}
MEMO-Appflow 是面向 Android 的端侧、基于本地大模型推理的应用预测与系统调度软件。它从 app 与操作系统两侧获取信息：在手机内持续采集真实 app 序列、eBPF 内核证据和系统状态，每 3 分钟压缩成 MAPLE 场景输入。大模型据此预测接下来可能使用的 Top-3 应用，并用 bitmask 选择调度动作，形成“从 app + OS 到 app + OS”的闭环：既更新桌面 Widget、优化 app 推荐体验，也执行预热、内存整理和服务刷新，改善 OS 状态。最新 9 分钟真机实验显示，Top-3 会随窗口刷新，\texttt{scheduler\_bits} 已触发 warm launch、memory trim 和 service manager 调度。
\end{abstract}

\tableofcontents
\newpage

\section{实时任务队列设计}

MEMO-Appflow 的实时模式不是“收完一段数据后停止等待推理”，而是把手机端采集和 MAPLE 推理解耦成一个有界流水线。系统持续运行一个采集流，每 3 分钟生成一个窗口。窗口结束后，当前窗口被压缩为 MAPLE scenario，并提交给唯一的 MAPLE worker；与此同时，下一个 3 分钟窗口立即继续采集。这样用户继续使用手机时，采集不会因为上一轮模型推理而暂停。

为了避免长期后台运行时线程或任务无限堆积，实时模式只保留三类运行单元：一个 collector loop、一个 MAPLE worker、一个 latest pending slot。如果 MAPLE 推理慢于窗口节奏，系统不会给每个窗口新开线程，也不会无限排队历史窗口，而是只保留最新待推理窗口。新窗口仍然以自己的 3 分钟 app/eBPF/system evidence 为主输入；EMA memory 只作为有界历史辅助信息，不替代当前窗口。

\begin{figure}[htbp]
\centering
\resizebox{\textwidth}{!}{
\begin{tikzpicture}[
  font=\small,
  node distance=0.9cm and 0.9cm,
  win/.style={draw, rounded corners, fill=blue!8, minimum width=3.2cm, minimum height=1.0cm, align=center},
  proc/.style={draw, rounded corners, fill=green!9, minimum width=3.1cm, minimum height=0.9cm, align=center},
  model/.style={draw, rounded corners, fill=orange!14, minimum width=3.2cm, minimum height=0.9cm, align=center},
  action/.style={draw, rounded corners, fill=purple!10, minimum width=3.2cm, minimum height=0.9cm, align=center},
  arrow/.style={-{Latex[length=2.2mm]}, thick}
]
\node[win] (w1) {窗口 1\\0--3 min\\采集 app + eBPF + OS};
\node[win, right=of w1] (w2) {窗口 2\\3--6 min\\继续采集};
\node[win, right=of w2] (w3) {窗口 3\\6--9 min\\继续采集};

\node[proc, below=of w1] (c1) {压缩窗口 1\\构造 scenario};
\node[proc, below=of w2] (c2) {压缩窗口 2\\构造 scenario};
\node[proc, below=of w3] (c3) {压缩窗口 3\\构造 scenario};

\node[model, below=of c1, xshift=1.55cm] (m1) {MAPLE worker\\推理窗口 1};
\node[model, below=of c2, xshift=1.55cm] (m2) {MAPLE worker\\推理窗口 2};
\node[model, below=of c3] (m3) {MAPLE worker\\推理窗口 3};

\node[action, below=of m1] (a1) {Top-3 + bitmask\\执行调度 1};
\node[action, below=of m2] (a2) {Top-3 + bitmask\\执行调度 2};
\node[action, below=of m3] (a3) {Top-3 + bitmask\\执行调度 3};

\draw[arrow] (w1) -- (w2);
\draw[arrow] (w2) -- (w3);
\draw[arrow] (w1) -- (c1);
\draw[arrow] (w2) -- (c2);
\draw[arrow] (w3) -- (c3);
\draw[arrow] (c1) -- (m1);
\draw[arrow] (c2) -- (m2);
\draw[arrow] (c3) -- (m3);
\draw[arrow] (m1) -- (a1);
\draw[arrow] (m2) -- (a2);
\draw[arrow] (m3) -- (a3);

\draw[arrow, dashed] ($(m1.north)+(0.7,0.1)$) -- node[above, sloped]{推理时窗口 2 不停采集} ($(w2.south)+(-0.4,-0.1)$);
\draw[arrow, dashed] ($(m2.north)+(0.7,0.1)$) -- node[above, sloped]{推理时窗口 3 不停采集} ($(w3.south)+(-0.4,-0.1)$);

\node[draw, dashed, rounded corners, fit=(m1)(m2)(m3), inner sep=0.25cm, label={[font=\small]below:单 MAPLE worker + latest pending slot，避免线程和任务爆炸}] {};
\end{tikzpicture}
}
\caption{3 分钟窗口采集与 MAPLE 异步推理流水线}
\end{figure}

\begin{table}[htbp]
\centering
\caption{实时任务队列的关键约束}
\begin{tabular}{p{0.25\textwidth}p{0.65\textwidth}}
\toprule
设计点 & 作用 \\
\midrule
3 分钟窗口 & 固定节奏聚合真实 app 序列、eBPF 事件和系统状态，避免每个事件都触发推理。\\
单 MAPLE worker & 控制端侧推理并发，避免本地小模型占用过多 CPU 和内存。\\
latest pending slot & 当 MAPLE 慢于采集节奏时只保留最新窗口，避免历史任务无限排队。\\
EMA memory & 保留有界历史偏好，但当前窗口仍然是主输入。\\
ActionExecutor & 把 Top-3 与 \texttt{scheduler\_bits} 转化为受控系统动作。\\
\bottomrule
\end{tabular}
\end{table}

\newpage
\section{引言与问题定义}

移动设备的性能体验来自两条线索：一条是用户接下来会打开什么 app，另一条是操作系统此刻是否有足够资源承接这个行为。传统应用预测通常只把问题建模成 app 序列预测，例如根据最近打开过的应用推荐下一个应用。这种方法可以得到一个 Top-3 列表，但它看不到更深层的系统状态：网络是否活跃、Binder 服务是否频繁、内存是否正在回收、显示线程是否繁忙、媒体解码是否正在发生。对于一个真正需要优化手机体验的系统，仅有 app 序列是不够的。

MEMO-Appflow 的核心思路是把 app 侧和 OS 侧信息放在同一条端侧闭环中。输入端同时收集真实前台 app 序列、eBPF 内核事件和 Android 系统状态；推理端用本地 MAPLE 模型预测用户接下来可能使用的 Top-3 真实应用，并判断系统应该执行哪些调度动作；输出端既更新桌面 Widget，让推荐成为可点击入口，也通过 ActionExecutor 执行系统动作，影响内存、预热、服务刷新和调度策略。概括来说，系统完成的是“从 app + OS 到 app + OS”：用 app 行为和 OS 证据共同赋能 app 推荐和 OS 性能优化。

这套系统的工程目标不是在电脑上跑一组脚本，而是让功能尽可能在 Android 设备内部闭环运行。Host 可以负责安装 APK、推送模型文件和拉取日志，但不能承担采集、解析、推理或动作执行的产品逻辑。这样设计的原因很直接：如果系统需要一直连着电脑才能工作，它就不是可以迁移到真实手机上的产品。

\section{端侧产品架构}

\subsection{闭环链路}

当前 MEMO-Appflow 的端侧链路如下：

\begin{verbatim}
真实前台 app 使用
  + raw eBPF 内核事件
  + Android 系统状态
        |
        v
3 分钟实时窗口聚合与压缩
        |
        v
MAPLE scenario JSON
        |
        v
设备本地 MAPLE 推理
        |
        v
Top-3 app id/category + scheduler_bits
        |
        v
真实 Android app 映射 + ActionExecutor
        |
        v
桌面 Widget 更新 + 系统调度动作
\end{verbatim}

这条链路中，用户看到的是 Top-3 真实应用和系统优化状态；eBPF、Binder、SurfaceFlinger、进程调度等信息不会直接堆在普通界面上，而是保存在报告和 debug 入口中。这样既保留了实验可追溯性，也避免把系统内部概念暴露成用户需要理解的产品功能。

\subsection{实时窗口与任务队列}

实时模式使用固定 3 分钟窗口。一个窗口结束后，系统将该窗口里的 app 序列、eBPF 事件和系统状态压缩成 scenario，并提交给唯一的 MAPLE worker；与此同时，下一个窗口继续采集。推理和采集解耦后，MAPLE 的 20-40 秒端侧推理耗时不会阻塞用户继续使用手机，也不会造成“收完一段、停下来推理、再继续收”的断点。

为了避免长期后台运行时线程和任务堆积，系统只保留一个 collector loop、一个 MAPLE worker 和一个 latest pending slot。如果 MAPLE 推理慢于窗口节奏，系统不会无限排队旧窗口，而是保留最新待推理窗口。当前窗口仍然是主输入；EMA memory 只保存有界历史偏好，作为辅助上下文，而不是无限增长的日志。

\begin{figure}[htbp]
\centering
\resizebox{\textwidth}{!}{
\begin{tikzpicture}[
  font=\small,
  node distance=1.0cm and 1.2cm,
  win/.style={draw, rounded corners, fill=blue!8, minimum width=3.3cm, minimum height=1.05cm, align=center},
  proc/.style={draw, rounded corners, fill=green!9, minimum width=4.2cm, minimum height=0.95cm, align=center},
  model/.style={draw, rounded corners, fill=orange!14, minimum width=4.6cm, minimum height=0.95cm, align=center},
  action/.style={draw, rounded corners, fill=purple!10, minimum width=4.6cm, minimum height=0.95cm, align=center},
  arrow/.style={-{Latex[length=2.2mm]}, thick}
]
\node[win] (w1) {窗口 1\\0--3 min\\持续采集};
\node[win, right=of w1] (w2) {窗口 2\\3--6 min\\持续采集};
\node[win, right=of w2] (w3) {窗口 3\\6--9 min\\持续采集};

\node[proc, below=1.2cm of w2] (compress) {窗口结束后压缩\\app timeline + eBPF counters + system state};
\node[model, below=0.85cm of compress] (worker) {单 MAPLE worker\\异步推理 latest pending window};
\node[action, below=0.85cm of worker] (out) {输出 Top-3 + scheduler bits\\更新 Widget 并执行系统动作};

\draw[arrow] (w1) -- (w2);
\draw[arrow] (w2) -- (w3);
\draw[arrow] (w1.south) |- (compress.west);
\draw[arrow] (w2) -- (compress);
\draw[arrow] (w3.south) |- (compress.east);
\draw[arrow] (compress) -- (worker);
\draw[arrow] (worker) -- (out);

\node[draw, dashed, rounded corners, fit=(compress)(worker)(out), inner sep=0.25cm] (runtime) {};
\node[below=0.15cm of runtime] {推理在后台执行；下一窗口继续采集；只保留单 worker 和最新待推理窗口};
\end{tikzpicture}
}
\caption{3 分钟采集窗口与 MAPLE 异步推理流水线}
\end{figure}

\section{数据结构与压缩}

\subsection{采集对象}

MAPLE 的输入不是单一 eBPF 日志，而是一个对齐后的 context JSON。它主要包含三类信息。

\textbf{App 侧信息}记录前台 app 序列，包括包名、应用名、切换时间和持续时间。它回答“刚才真实用了哪些 app”。这部分对于 Top-3 应用预测是必要的，因为用户行为本身有连续性。

\textbf{OS 侧信息}来自 eBPF 和系统状态采样。eBPF 事件包括网络 sendto/recvfrom、文件 openat、Binder/服务调用、进程生命周期、调度和运行时信号。系统状态包括 memory、battery、network、camera/media、display/UI、process/service 等。它回答“这些 app 使用发生时，系统内部处在什么状态”。

\textbf{可执行动作契约}列出 Android 端已经实现的调度动作。MAPLE 不直接输出 shell 命令，而是在这组动作中用 0/1 bitmask 选择是否执行，降低模型直接控制系统的风险。

\subsection{时间对齐与重复压缩}

仅有 app 序列或仅有 eBPF 都不够。一个合理输入应该告诉模型：同一个时间段里，哪个 app 在前台，OS 内部发生了哪些事件，事件强度是多少。因此系统按时间窗口把 app segment 和 eBPF counter 对齐。比如某个窗口中 Chrome 在前台，同时出现大量 recvfrom 和 display refresh，模型就能理解这是浏览/媒体/网络相关场景，而不是孤立地看到一串网络系统调用。

原始 eBPF trace 会有大量重复行。如果把重复事件逐条送进 prompt，模型上下文会被低信息密度文本占满。MEMO 的压缩方式是把同一个窗口内相同类型、相同证据类别、相近进程来源的事件合并成 counter key，并保留 count、时间范围和代表性样本。例如上千条 \texttt{MEMO\_RECVFROM} 会压缩成“Network IO 在该窗口重复 N 次，主要关联哪些进程和前台 app”。这样会损失单条事件细节，但保留调度最需要的信息：发生了什么、发生在哪个窗口、强度有多高、和哪个 app 使用阶段相关。

\section{MAPLE 推理与系统动作}

\subsection{MAPLE 输入输出}

MAPLE 是唯一预测入口。Android 端构造 scenario JSON 后，本地调用 MAPLE/llama.cpp 可执行文件和 GGUF 模型。Prompt 明确要求预测“用户接下来可能使用的三个 app”，同时输出系统调度 bitmask。最终输出分为三层：

\begin{itemize}
  \item Stage 1 预测用户可见 app 类别，例如 Communication、Media、Browser、Camera。
  \item Stage 2 在类别候选中选择 MAPLE app id，再由 Android 映射成真实包名。
  \item Stage 3 输出 \texttt{scheduler\_bits}，选择系统动作。
\end{itemize}

这里需要区分 app prediction candidates 和 scheduler evidence。Memory Management、Binder、process runtime、system service 这类 OS-only 类别可以影响调度，但不能被当成“下一个 app 类别”。否则小模型容易把系统资源类别误预测为用户要点击的应用类别。当前版本将用户可见 app 类别和 OS 调度证据分开：前者用于 Top-3，后者用于 \texttt{scheduler\_bits} 和安全策略。

\subsection{真实应用映射与 Widget}

MAPLE 输出的 App id 或类别不能直接展示给用户。AppIdMapping 使用 PackageManager 扫描设备上的 launchable apps，读取真实包名、label、launch intent 和可用类别，再根据 MAPLE 类别、历史使用、包名/label 关键词和设备实际安装情况排序。这样迁移到另一台手机时不需要硬编码“必须有某个 app”：如果通信类场景中没有 QQ，但有微信、短信或电话，Communication 类候选会自然落到那台设备上真实存在且可启动的应用。

桌面 Widget 是产品出口。实时模式每 3 分钟刷新 Top-3 后，Widget 展示真实 app 图标、名称和更新时间，点击图标可以直接打开应用。App 内部则展示更详细的系统判断、证据摘要、动作日志和报告入口。

\subsection{ActionExecutor 与 scheduler bits}

ActionExecutor 把 MAPLE 的 bitmask 转成系统动作。当前定义 9 个动作位。表中使用动作简称；代码中对应的完整 action id 保存在 `ActionExecutor.kt` 的动作表中。

\begin{longtable}{r p{0.26\textwidth} p{0.54\textwidth}}
\caption{scheduler bits 对应的 Android 端动作}\\
\toprule
bit & 动作简称 & 含义 \\
\midrule
\endfirsthead
\toprule
bit & 动作简称 & 含义 \\
\midrule
\endhead
0 & \texttt{warm\_top1} & 预热 Top-1 推荐应用并回到 HOME。\\
1 & \texttt{warm\_top2\_idle} & 设备空闲时预热 Top-2。\\
2 & \texttt{memory\_trim} & 对低优先级候选应用做 memory trim。\\
3 & \texttt{kill\_background} & 关闭被选中的低优先级后台包。\\
4 & \texttt{drop\_cache} & 仅在 critical memory 压力下 drop cache。\\
5 & \texttt{network\_refresh} & 刷新网络统计。\\
6 & \texttt{service\_refresh} & 刷新或查询 service manager。\\
7 & \texttt{display\_guard} & 显示/UI 繁忙时降低预热强度。\\
8 & \texttt{thermal\_guard} & 温度高时跳过 camera/media 预热。\\
\bottomrule
\end{longtable}

模型负责选择策略，Android 端负责安全执行。比如 \texttt{scheduler\_bits=101010100} 会请求 bit0、bit2、bit4、bit6 等动作，但 drop cache 仍然需要通过安全检查；如果当前内存不是 critical，ActionExecutor 会记录 skipped，而不是盲目执行。

\section{实验设计}

本文只讨论最终成功版本对应的实验。实验回答三个问题：

\begin{itemize}
  \item 实时产品闭环是否能运行：连续采集、连续推理、刷新 Top-3 和执行调度。
  \item eBPF/OS 信息是否有助于预测：与只看 app 序列的 baseline 对比。
  \item Stage 3 调度是否真的影响系统指标：同一 Top-1 app 上比较无调度和有调度。
\end{itemize}

\textbf{9 分钟实时 Top-3 实验}在 rooted Pixel 5 上运行三个 3 分钟窗口。每个窗口真实启动设备上存在的 app，并收集自然产生的 eBPF/system evidence。每个窗口结束后由 MAPLE 推理并刷新 Top-3 和调度动作。该实验验证产品形态。

\textbf{受控理论消融实验}使用 Chengyu Fan 构造的 1000 条带 ground truth 的合成样本，逐组删除 app sequence、Binder、file、memory 等证据，观察多源 eBPF/OS 证据是否比 app sequence only 更有信息量。它先在可控数据上验证机制。

\textbf{30 条真实记录预测实验}使用 3 个 persona，每个 persona 10 条短 app 使用记录，共 30 条。Codex 根据 persona 生成“已观察 app 序列”和被隐藏的未来 Top-3 ground truth；执行阶段在 Android 设备上真实启动 app、发送 input 命令，并自然产生 raw eBPF。评估时模型看不到未来 Top-3，只看到输入 app 序列和压缩 OS/eBPF context。

\textbf{系统调度性能实验}不评价预测准确率，只评价 Stage 3 动作是否改善系统指标。实验先用真实 eBPF + MAPLE 选出同一个 Top-1 目标 app，本次为 QQ。Baseline 分支 force-stop QQ 后直接启动并测量；Stage 3 分支 force-stop QQ 后先执行同一组调度 bitmask，再启动同一个 QQ。这样比较对象一致，差别集中在是否执行 MAPLE 调度。

\section{结果与分析}

\subsection{实时闭环结果}

9 分钟实验的核心结果见表 \ref{tab:realtime}。三个窗口均产生约 1.9 万条 raw eBPF 事件和 52-55 个 app segment，说明实验不是手写 JSON，而是在手机上真实产生应用切换和系统事件。Top-3 在窗口间变化两次，说明推荐不是固定列表，而是随最近窗口更新。

\begin{table}[htbp]
\centering
\caption{9 分钟实时 Top-3 实验结果}
\label{tab:realtime}
\resizebox{\textwidth}{!}{
\begin{tabular}{lrrrrl}
\toprule
窗口 & raw eBPF & app segments & MAPLE time & App id & Top-3 \\
\midrule
network & 19263 & 52 & 29106 ms & 220 & Messages / Camera / Chrome \\
camera & 19289 & 55 & 37196 ms & 115 & QQ / Messages / Camera \\
media\_comm & 19260 & 54 & 23839 ms & 220 & Messages / Camera / QQ \\
\bottomrule
\end{tabular}
}
\end{table}

调度方面，前两个窗口输出 \texttt{100000000}，执行 Top-1 warm launch；第三个窗口输出 \texttt{101010100}，执行 widget update、warm launch、memory trim 和 service manager refresh，并因为内存不是 critical 安全跳过 drop cache。这说明 MAPLE 输出已经进入 Android 端动作执行层，而不是停留在文本建议。

\subsection{预测消融结果}

受控理论消融先验证 eBPF/OS 证据本身有用。1000 条合成样本中，Full eBPF 在 Qwen CPU 上达到 63.0\%，App sequence only 为 30.1\%，Random 为 4.8\%。DeepSeek 配置下 Full eBPF 为 88.8\%，App sequence only 为 51.1\%。这说明多源 OS 证据在可控 ground truth 上提供了明显额外信息。

真实设备的 30 条记录实验得到同方向结果。表 \ref{tab:prediction} 展示最终成功配置：MEMO app+OS/eBPF 在 Top-1 Accuracy、Top-3 Hit、Precision@3、Recall@3、F1@3 和 MRR 上都高于 app-only baseline。提升幅度不夸张，但方向与理论消融一致，说明 eBPF/OS context 在真实设备路径中也能提供增益。

\begin{table}[htbp]
\centering
\caption{30 条真实记录预测实验结果}
\label{tab:prediction}
\resizebox{\textwidth}{!}{
\begin{tabular}{lrrrrrr}
\toprule
配置 & Top-1 Acc & Top-3 Hit & Precision@3 & Recall@3 & F1@3 & MRR \\
\midrule
App-only baseline & 20.00\% & 46.67\% & 47.78\% & 47.78\% & 47.78\% & 31.11\% \\
\textbf{MEMO app+OS/eBPF} & \textbf{23.33\%} & \textbf{50.00\%} & \textbf{52.22\%} & \textbf{52.22\%} & \textbf{52.22\%} & \textbf{34.44\%} \\
\bottomrule
\end{tabular}
}
\end{table}

相对于 app-only baseline，Top-3 Hit 提升 3.33 个百分点，F1@3 提升 4.44 个百分点，Top-1 Accuracy 提升 3.33 个百分点，MRR 提升 3.33 个百分点。这里的结论应谨慎表述：eBPF/OS 并不是单独替代 app 序列，而是作为系统上下文补充 app 行为预测，并为后续调度提供证据。

\subsection{系统调度性能结果}

系统调度性能实验验证 Stage 3 是否真的改变系统指标。结果见表 \ref{tab:scheduler}。在同一 MAPLE Top-1 app QQ 上，Stage 3 调度后 Launch TotalTime 从 433 ms 降到 168 ms，WaitTime 从 440 ms 降到 173 ms，综合压力分数从 422.56 降到 9.71。CPU busy、iowait、reclaim 和 jank rate 也下降。

\begin{table}[htbp]
\centering
\caption{系统调度性能实验结果}
\label{tab:scheduler}
\begin{tabular}{lrrr}
\toprule
指标 & No scheduler & Stage 3 & 改善 \\
\midrule
Launch TotalTime & 433 ms & 168 ms & 61.20\% \\
WaitTime & 440 ms & 173 ms & 60.68\% \\
综合压力分数 & 422.56 & 9.71 & 97.70\% \\
CPU busy & 26.70\% & 20.79\% & 22.13\% \\
iowait & 0.094\% & 0.023\% & 75.42\% \\
MemAvailable drop & 164892 KB & -84448 KB & 151.21\% \\
reclaim & 10046 & 0 & 100.00\% \\
jank rate & 1.240\% & 0.886\% & 28.56\% \\
\bottomrule
\end{tabular}
\end{table}

这组结果的关键不在于声称所有场景都会提升，而在于证明 MAPLE 调度 bitmask 已经可以进入 ActionExecutor，触发系统动作，并在可测指标上留下结果。Stage 3 实际执行了 7 条动作记录，其中包括 Top-1 warm launch、memory trim 和 service manager refresh；drop cache 被安全规则跳过，因为当时内存不是 critical。

\section{实现对应关系}

\begin{longtable}{p{0.26\textwidth} p{0.64\textwidth}}
\caption{关键实现文件}\\
\toprule
模块 & 作用 \\
\midrule
\endfirsthead
\toprule
模块 & 作用 \\
\midrule
\endhead
实时总控 & \texttt{RealtimePreloadController.kt}：管理持续采集、窗口提交、MAPLE worker 和调度执行。\\
窗口采集 & \texttt{RealtimeWindowRunner.kt}：负责每个窗口内的 app/eBPF/system evidence 收集和压缩。\\
历史记忆 & \texttt{RealtimeMemoryStore.kt}：维护有界 EMA memory，避免历史上下文无限增长。\\
场景构造 & \texttt{MapleScenarioBuilder.kt}：构造 MAPLE scenario JSON，并分离 app prediction candidates 与 scheduler evidence。\\
模型调用 & \texttt{MapleShellBackend.kt}：在 Android 端调用本地 MAPLE/llama.cpp 可执行文件并解析输出。\\
应用映射 & \texttt{AppIdMapping.kt}：扫描真实安装应用，把 MAPLE id/category 映射成可启动 Top-3 app。\\
动作执行 & \texttt{ActionExecutor.kt}：执行调度 bitmask 对应的系统动作，并记录执行/跳过原因。\\
桌面入口 & \texttt{MemoWidgetProvider.kt}：更新桌面 Widget，展示可点击 Top-3 真实应用和最新更新时间。\\
性能实验 & \texttt{PressureExperimentRunner.kt}：执行系统调度性能实验，比较同一目标 app 的无调度与 Stage 3 调度。\\
\bottomrule
\end{longtable}

\section{局限与下一步}

当前系统仍有局限。第一，本地小模型不够“听话”，有时输出格式不完全符合 prompt。项目已经增加 parser 处理脏输出，但还不够 robust，后续应继续加强结构化输出约束和异常恢复。第二，实时性和轻量性存在 tradeoff。每 3 分钟推理一次可以让推荐和调度更及时，但后台持续采集和本地推理会消耗 CPU、电量和内存。第三，app scanner 对部分类别的映射还可以更精细，例如 camera 类场景应更稳定地选到真实相机应用。

下一步可以沿三个方向推进：其一，将 MAPLE 输出协议进一步结构化，减少自由文本解析；其二，增强 AppIdMapping 的自动分类能力，让不同手机上的真实应用映射更准确；其三，把调度动作从当前 warm launch、memory trim、service refresh 扩展到更细粒度的 Android 资源策略，同时继续保留安全门控，避免模型输出直接触发高风险系统操作。

\section{结论}

MEMO-Appflow 已经从 host 脚本原型推进为 Android 端侧产品闭环。它把 app 序列和 OS/eBPF 证据按时间对齐并压缩，交给本地 MAPLE 模型推理，再将结果转成真实 Top-3 app、桌面 Widget 和系统调度动作。实验表明，实时 Top-3 能随窗口更新；app+OS/eBPF 在真实记录预测中优于 app-only baseline；Stage 3 调度能在同一目标 app 上显著降低启动延迟和系统压力指标。更重要的是，这些结果来自手机内部运行的采集、推理和执行链路，而不是离线修改 JSON 或 host 侧脚本模拟。


\end{document}
